%%%%%%%%%%%%%%%%%%%%%%%%%%%%%%%%%%%%%%%%%%%%%%%%%%%%%%%%%%%%%%%%%%%%%%%%%%%
%
% Generic template for TFC/TFM/TFG/Tesis
%
% $Id: desarrollo.tex,v 1.4 2015/06/05 00:05:18 macias Exp $
%
% By:
%  + Javier Macías-Guarasa.
%    Departamento de Electrónica
%    Universidad de Alcalá
%  + Roberto Barra-Chicote.
%    Departamento de Ingeniería Electrónica
%    Universidad Politécnica de Madrid
% 
% Based on original sources by Roberto Barra, Manuel Ocaña, Jesús Nuevo, Pedro Revenga, Fernando Herránz and Noelia Hernández. Thanks a lot to all of them, and to the many anonymous contributors found (thanks to google) that provided help in setting all this up.
%
% See also the additionalContributors.txt file to check the name of additional contributors to this work.
%
% If you think you can add pieces of relevant/useful examples, improvements, please contact us at (macias@depeca.uah.es)
%
% You can freely use this template and please contribute with comments or suggestions!!!
%
%%%%%%%%%%%%%%%%%%%%%%%%%%%%%%%%%%%%%%%%%%%%%%%%%%%%%%%%%%%%%%%%%%%%%%%%%%%

\chapter{Desarrollo}
\label{cha:desarrollo}

\begin{FraseCelebre}
  \begin{Frase}
    A fuerza de construir bien, se llega a buen arquitecto\footnote{Tomado de ejemplos del proyecto \texis{}.}.
  \end{Frase}
  \begin{Fuente}
    Aristóteles
  \end{Fuente}
\end{FraseCelebre}


\section{Introducción/Metodología}
\label{sec:introduccion-desarrollo}

En este capítulo se incluirá la descripción del desarrollo del trabajo.

El capítulo se estructura en n apartados....


\section{Dataset empleado}
\label{sec:desarr-del-sist}

La elección del dataset constituye una parte fundamental en el desarrollo de este trabajo, ya que condiciona tanto la metodología empleada como la validez de los resultados obtenidos. Inicialmente, el objetivo del TFG se centraba en la detección de ansiedad y depresión a partir de texto, pero pronto se observó que encontrar corpus adecuados para este propósito resultaba especialmente complejo. La mayoría de datasets públicos disponibles incluían únicamente casos de depresión, ofrecían información limitada sobre la calidad del etiquetado o no contaban con un proceso de anotación suficientemente fiable.

Tras estudiar distintas alternativas, se seleccionó finalmente el dataset \textit{MentalRiskES}, creado por investigadores de la Universidad de Jaén junto con colaboradores internacionales. Este corpus proporciona anotaciones realizadas por múltiples evaluadores y presenta un mayor grado de consistencia que otros datasets disponibles. Además, MentalRiskES no solo incluye datos sobre ansiedad y depresión, sino también sobre trastornos de la conducta alimentaria, por lo que se decidió ampliar el alcance del trabajo para incorporar la detección de este tercer trastorno. 

Un aspecto especialmente relevante es que MentalRiskES está compuesto íntegramente por texto en español, lo que permite desarrollar un estudio que se aleja de la tendencia dominante en la literatura, centrada mayoritariamente en corpus en lengua inglesa.


\subsection{Análisis del dataset MentalRiskES}

\subsubsection{Descripción y estructura general del corpus}
MentalRiskES constituye un corpus en español diseñado para la detección temprana de tres trastornos mentales: depresión, ansiedad y trastornos de la conducta alimentaria (Eating Disorders, ED). Los datos que componen el corpus provienen de grupos públicos de Telegram, seleccionados específicamente por su naturaleza temática relacionada con cada uno de los trastornos.

La unidad básica del dataset es el sujeto (usuario). Para cada usuario se incluye un identificador anonimizado y una secuencia cronológica de mensajes publicados dentro de un grupo concreto. El dataset conserva los mensajes originales tras un proceso de curación que incluyó la eliminación de mensajes extremadamente breves, la conversión de emojis a texto, el etiquetado de URLs y hashtags y la anonimización de nombres propios, alias y números de teléfono.

El dataset está organizado en tres carpetas, una por cada trastorno. En cada una se encuentra un archivo \texttt{JSON} por sujeto, que contiene todos sus mensajes ordenados cronológicamente, junto con un archivo \texttt{CSV} que recoge el identificador del sujeto y dos etiquetas:
\begin{itemize}
	\item \textbf{bs (binary suffer):} indica si los mensajes publicados por el sujeto presentan señales de sufrir el trastorno (1) o no (0).
	\item \textbf{bc (binary control):} indica si el sujeto debe considerarse control (1) o no (0). Se entiende por control a aquellos usuarios que no muestran señales asociadas al trastorno en cuestión.
\end{itemize}

\subsubsection{Proceso de anotación}
Es importante destacar que estas etiquetas no fueron establecidas por profesionales de la salud mental. El proceso de anotación fue realizado por diez anotadores independientes por sujeto, quienes, aunque no eran especialistas clínicos, cumplían criterios estrictos de selección (ser nativos en español y poseer estudios universitarios) y trabajaron siguiendo guías detalladas basadas en definiciones oficiales y procedimientos estandarizados. La etiqueta final asignada a cada sujeto se determina seleccionando la que aparece con mayor frecuencia entre los diez anotadores y, en caso de empate, se considera que el usuario sí presenta indicios del trastorno correspondiente.

Este enfoque proporciona un nivel de fiabilidad adecuado en las anotaciones obtenidas. No obstante, la ausencia de anotadores profesionales implica que las etiquetas deben interpretarse como indicadores de riesgo y no como diagnósticos clínicos. Por ello, futuros trabajos podrían beneficiarse de corpus anotados por especialistas en salud mental para aumentar aún más su fiabilidad.

\subsubsection{Exploración del dataset}
A continuación, se presenta una exploración inicial del corpus \textit{MentalRiskES} con el objetivo de entender su distribución de usuarios, mensajes y palabras, así como el balance entre clases para cada trastorno. Para ello se generó un pequeño script en Python que permitió extraer los siguientes resultados:

\begin{itemize}
	\item En los trastornos de la conducta alimentaria (ED), el corpus contiene 335 sujetos, de los cuales 192 pertenecen a la clase control y 143 a la clase \textit{suffer}. Aunque existe cierto desbalance, la diferencia entre clases es moderada.
	\item En depresión, la proporción está más desequilibrada. De los 499 sujetos presentes en esta partición, 166 son control y 333 se agrupan dentro de las distintas subcategorías de \textit{suffer}, lo que supone aproximadamente el doble de sujetos positivos respecto a negativos.
	\item La partición correspondiente a ansiedad presenta el desbalance más pronunciado: solo 57 sujetos pertenecen a la clase control, frente a 443 que han sido anotados como \textit{suffer}. Este desequilibrio debe tenerse en cuenta en fases posteriores, ya que puede afectar al rendimiento y a la sensibilidad del modelo.
\end{itemize}

Para visualizar estas diferencias, la Figura~\ref{fig:dist_clases_mentalriskes} muestra un diagrama de barras agrupadas con la distribución de clases en los tres trastornos del corpus.

\begin{figure}[H]
	\centering
	\includegraphics[width=0.8\textwidth]{distribucion_clases_mentalriskes.png}
	\caption{Distribución de clases en los tres trastornos del corpus \textit{MentalRiskES}.}
	\label{fig:dist_clases_mentalriskes}
\end{figure}

Otro factor importante a tener en cuenta es el número de mensajes y de palabras que aporta cada usuario, ya que determina la cantidad de información textual disponible para los modelos de clasificación. Para analizar esta cuestión, se calculó, para cada trastorno y clase, el número total de mensajes y palabras, así como sus medias por usuario. Los resultados se resumen en la Tabla~\ref{tab:resumen_trastornos}.

\begin{table}[H]
	\centering
	\caption{Resumen de estadísticas por trastorno y clase en el corpus MentalRiskES.}
	\renewcommand{\arraystretch}{1.2}
	\setlength{\tabcolsep}{6pt}
	\resizebox{\textwidth}{!}{
		\begin{tabular}{|l|ccc|ccc|ccc|}
			\hline
			& \multicolumn{3}{c|}{\textbf{Ansiedad}} 
			& \multicolumn{3}{c|}{\textbf{Depresión}} 
			& \multicolumn{3}{c|}{\textbf{ED}} \\
			\textbf{Métrica} 
			& \textbf{Control} & \textbf{Suffer} & \textbf{Total}
			& \textbf{Control} & \textbf{Suffer} & \textbf{Total}
			& \textbf{Control} & \textbf{Suffer} & \textbf{Total} \\
			\hline
			\textbf{Nº sujetos}      
			& 57  & 443 & 500
			& 166 & 333 & 499
			& 192 & 143 & 335 \\
			\hline
			\textbf{Nº mensajes}     
			& 2\,222  & 16\,293 & 18\,515
			& 6\,287  & 11\,083 & 17\,370
			& 6\,586  & 3\,913  & 10\,499 \\
			\textbf{Media mensajes por usuario}   
			& 38.98 & 36.78 & 37.03
			& 37.87 & 33.28 & 34.81
			& 34.30 & 27.36 & 31.34 \\
			\hline
			\textbf{Nº palabras}     
			& 31\,248  & 276\,853 & 308\,101
			& 84\,435  & 159\,868 & 244\,303
			& 93\,614  & 64\,667  & 158\,281 \\
			\textbf{Media palabras por usuario}   
			& 548.21 & 624.95 & 616.20
			& 508.64 & 480.08 & 489.59
			& 487.57 & 452.22 & 472.48 \\
			\hline
		\end{tabular}
	}
	\label{tab:resumen_trastornos}
\end{table}

El análisis del número de mensajes y palabras revela diferencias relevantes entre los tres trastornos. Ansiedad es el conjunto con mayor volumen total de mensajes y palabras, seguido de depresión y, finalmente, ED. Si se comparan las medias por usuario, en ansiedad y depresión los usuarios de la clase control tienden a escribir ligeramente más mensajes que los usuarios \textit{suffer}, mientras que en ED ocurre lo contrario. En cuanto al número de palabras, en ansiedad los usuarios \textit{suffer} producen mensajes más extensos de media, en depresión sucede lo contrario, y en ED ambas clases presentan valores muy similares.

\subsubsection{Conclusión del análisis del dataset}

En conjunto, los resultados muestran diferencias claras tanto en el desbalance entre clases como en la cantidad de texto disponible por trastorno. Ansiedad es el conjunto más desequilibrado, mientras que depresión y ED presentan distribuciones más homogéneas. Asimismo, ansiedad y depresión ofrecen un volumen textual mayor que ED, lo que implica que los modelos dispondrán de más información lingüística en estos dos casos. Estas variaciones son relevantes porque afectan tanto a la dificultad del aprendizaje como a la capacidad de generalización, por lo que deben tenerse en cuenta tanto en el proceso de entrenamiento del modelo BERT como en la evaluación posterior, tanto de BERT como de los modelos LLM preentrenados.


\subsection{Preparación de datos}

Una vez analizado el dataset, el siguiente paso consistió en preparar los datos para su uso en los modelos de clasificación. Dado que el dataset \textit{MentalRiskES} está estructurado a nivel de usuario y las etiquetas asignadas corresponden al sujeto en su conjunto, se decidió mantener esta unidad de análisis durante todo el proceso. Así, los mensajes asociados a cada usuario se concatenaron en un único texto, preservando la información contextual completa que aporta el corpus.

El proceso comenzó con la carga de los ficheros \texttt{JSON} de cada trastorno, a partir de los cuales se extrajo la lista de mensajes correspondiente a cada usuario. Estos mensajes fueron filtrados para eliminar entradas vacías y posteriormente concatenados en un único texto. Paralelamente, se incorporó la etiqueta binaria \texttt{bs} asociada a cada sujeto, generando tres \textit{dataframes} finales (uno por trastorno) con las columnas \texttt{user\_id}, \texttt{texto} y \texttt{bs}.

\subsubsection{Partición de datos para clasificación binaria}

En primer lugar, se construyeron los conjuntos destinados a la clasificación binaria, en los que el objetivo es determinar si un usuario presenta (\textit{suffer}) o no presenta (\textit{control}) un determinado trastorno. A partir de los dataframes individuales, se procedió a preparar los datos para el modelo BERT. Para este tipo de modelos es necesario dividir cada dataset en tres subconjuntos: entrenamiento (70\%), validación (15\%) y prueba (15\%). El conjunto de entrenamiento se emplea para ajustar los parámetros del modelo, el de validación para seleccionar hiperparámetros y prevenir el sobreajuste, y el de prueba para evaluar el rendimiento final.

Dado que, tal como se observó en el análisis del subapartado 3.2.1.3, las clases \textit{suffer} y \textit{control} presentan un desbalance significativo (especialmente en el caso del trastorno de ansiedad), la división se realizó mediante estratificación. Esto garantiza que la proporción entre ambas clases se conserve en los subconjuntos de entrenamiento, validación y prueba, manteniendo la distribución del corpus original. Los subconjuntos resultantes se almacenaron en ficheros CSV para su uso posterior en el proceso de entrenamiento de BERT.

En el caso de los modelos LLM, el procedimiento es más simple, ya que no requieren entrenamiento adicional. Por ello, no se generan subconjuntos de entrenamiento ni de validación, sino únicamente un conjunto de evaluación por trastorno. Dado que los conjuntos de prueba estratificados resultan demasiado pequeños y presentan un fuerte desbalance (especialmente en ansiedad), se decidió emplear el dataset completo de cada trastorno, ya con sus textos unificados, como conjunto de evaluación para los LLM.

\subsubsection{Partición de datos para clasificación multiclase}

Además del escenario binario, se prepararon también los datos para abordar una clasificación multiclase, en la que el objetivo consiste en identificar cuál de los tres trastornos presenta un usuario, o si pertenece a la clase control. Para ello se fusionaron los dataframes procesados de ansiedad, depresión y trastornos alimentarios en un único dataset. A cada usuario se le asignó una etiqueta multiclase en función de su trastorno de origen (ansiedad, depresión o ED) o del estado \textit{control} si su etiqueta binaria era 0. Este proceso dio lugar a cuatro clases en total: control, ansiedad, depresión y ED.

El dataset multiclase resultante se dividió posteriormente en los mismos subconjuntos que en el caso binario (70\% entrenamiento, 15\% validación y 15\% prueba), aplicando una división estratificada para preservar la proporción entre clases. Este conjunto se empleará para entrenar una variante multiclase de BERT.

De manera análoga al caso binario, los modelos LLM se evaluarán sobre el dataset multiclase completo, sin aplicar división en subconjuntos, ya que no requieren un proceso de entrenamiento previo.

\subsubsection{Subconjuntos para la métrica TARa@5}

Con el objetivo de calcular la métrica \ac{TARa@5}, utilizada para medir la estabilidad de las predicciones de los \acp{LLM} y descrita previamente en la Sección~\ref{sec:metricas_evaluacion}, se prepararon subconjuntos específicos de evaluación para cada tarea de clasificación.

Dado que calcular esta métrica sobre el dataset completo supondría un coste de inferencia muy elevado (ya que cada ejemplo debe evaluarse cinco veces por cada modelo), se optó por construir subconjuntos reducidos pero representativos mediante muestreo estratificado. De este modo, se conserva la proporción original entre clases y se garantiza que la muestra siga siendo representativa del problema de clasificación.

En los escenarios de clasificación binaria, se seleccionaron 100 ejemplos por trastorno a partir de los conjuntos completos utilizados para la inferencia con \acp{LLM}. Esta cantidad representa aproximadamente un 20\% del tamaño total de los datasets en los casos de ansiedad y depresión, y cerca de un 30\% en el caso de \ac{TCA}, porcentaje ligeramente mayor debido al menor tamaño del dataset original, ya que una proporción inferior habría generado un subconjunto demasiado reducido.

En el caso de la clasificación multiclase, se extrajo un subconjunto de 250 ejemplos manteniendo igualmente la distribución original entre las cuatro clases (control, ansiedad, depresión y ED). Esta cantidad supone aproximadamente un 20\% del dataset multiclase completo, proporcionando representación suficiente de cada categoría y una estimación estable de la consistencia del modelo en un escenario más complejo.

\subsection{Análisis complementarios previos a la clasificación}

Antes de iniciar el entrenamiento y la evaluación de los modelos, se realizaron varios análisis complementarios orientados a validar algunas decisiones metodológicas y comprender mejor la estructura del problema de clasificación. 

\subsubsection{Longitud de secuencia y justificación del uso de Longformer}

En primer lugar, se estudió la longitud de los textos generados tras la concatenación de los mensajes de cada usuario. Este paso resulta especialmente relevante porque los modelos \ac{BERT} estándar presentan una limitación estructural de 512 tokens de entrada, por lo que era necesario comprobar si los textos del corpus superaban dicho umbral.

Para ello, se utilizó el tokenizador base de \ac{BERT} y se calculó el número total de tokens de cada usuario en los distintos datasets preparados (ansiedad, depresión, \ac{TCA} y multiclase). La Tabla~\ref{tab:token_lengths} resume las principales estadísticas obtenidas.

\begin{table}[H]
	\centering
	\caption{Longitud media y máxima de tokens por usuario tras la concatenación de mensajes}
	\label{tab:token_lengths}
	\begin{tabular}{|l|c|c|}
		\hline
		\textbf{Categoría} & \textbf{Media de tokens} & \textbf{Máximo de tokens} \\
		\hline
		Ansiedad & 1042.98 & 6134 \\
		Depresión & 820.24 & 6457 \\
		\ac{TCA} & 799.53 & 12754 \\
		Multiclase & 898.52 & 12754 \\
		\hline
	\end{tabular}
\end{table}

Los resultados muestran que, en todos los casos, la longitud media supera ampliamente el límite de 512 tokens, alcanzando además valores máximos muy superiores. Este análisis permitió identificar una de las primeras dificultades encontradas durante el desarrollo del trabajo: una parte importante de la información contextual quedaría fuera si se utilizara directamente un modelo \ac{BERT} estándar con truncación, debido a su límite máximo de 512 tokens de entrada.

Ante esta situación, existen dos estrategias principales. La primera consiste en aplicar truncación durante la tokenización, conservando únicamente los primeros 512 tokens de cada texto. Esta solución resulta sencilla y eficiente, aunque puede provocar pérdida de información relevante. La segunda opción consiste en utilizar modelos Transformer diseñados para secuencias largas, capaces de procesar miles de tokens sin necesidad de truncar, aunque con un mayor coste computacional.

En este trabajo se decidió emplear ambos enfoques: por un lado, modelos \ac{BERT} convencionales con truncación a 512 tokens y, por otro, un modelo \textit{long-range}, concretamente Longformer, capaz de procesar secuencias mucho más extensas. Esto permite analizar hasta qué punto la truncación afecta al rendimiento de clasificación y si el mayor coste computacional de los modelos de largo contexto se traduce realmente en una mejora significativa en el rendimiento de clasificación.

\subsubsection{Visualización de embeddings y separabilidad entre clases}

En segundo lugar, se realizó un análisis exploratorio sobre la separabilidad semántica de las clases antes del entrenamiento de los modelos. El objetivo de este estudio fue comprobar si los textos asociados a ansiedad, depresión, \ac{TCA} y control presentaban de forma natural agrupaciones diferenciadas en el espacio vectorial o, por el contrario, existía un fuerte solapamiento entre categorías.

Para ello, se generaron embeddings de cada texto utilizando el modelo \texttt{sentence-transformers/paraphrase-multilingual-mpnet-base-v2} \cite{reimers_making_2020}, un modelo de \textit{sentence embeddings} multilingüe diseñado para representar documentos en forma de vectores densos que capturan su significado semántico. Cada texto del dataset multiclase se transformó así en una representación numérica de alta dimensión.

Dado que estos embeddings se representan originalmente en un espacio vectorial de alta dimensión ($\mathbb{R}^{768}$), su visualización directa no resulta posible. Por ello, se proyectaron a un espacio bidimensional ($\mathbb{R}^2$) mediante dos técnicas de reducción de dimensionalidad: \ac{PCA} \cite{jolliffe_principal_2016} y \ac{t-SNE} \cite{van_der_maaten_visualizing_2008}. Mientras que \ac{PCA} permite observar la estructura global de los datos preservando la mayor varianza posible, \ac{t-SNE} facilita una visualización más intuitiva de agrupaciones locales y relaciones de vecindad entre ejemplos.

Las Figuras~\ref{fig:embeddings_visualization} muestran las proyecciones obtenidas para el dataset multiclase, utilizando una codificación por color para cada clase: control (gris), ansiedad (rojo), depresión (azul) y \ac{TCA} (verde).

\begin{figure}[H]
	\centering
	
	\begin{subfigure}[b]{0.49\textwidth}
		\centering
		\includegraphics[width=\textwidth]{embeddings_pca.png}
		\caption{PCA}
		\label{fig:embeddings_pca}
	\end{subfigure}
	\hfill
	\begin{subfigure}[b]{0.49\textwidth}
		\centering
		\includegraphics[width=\textwidth]{embeddings_tsne.png}
		\caption{t-SNE}
		\label{fig:embeddings_tsne}
	\end{subfigure}
	
	\caption{Visualización de embeddings del dataset multiclase mediante reducción de dimensionalidad}
	\label{fig:embeddings_visualization}
\end{figure}

En la proyección mediante \ac{PCA} se observa una estructura relativamente organizada, donde la clase \ac{TCA} aparece más concentrada y desplazada hacia una región más diferenciada del espacio vectorial, lo que sugiere que sus patrones lingüísticos presentan mayor especificidad semántica. Por el contrario, ansiedad y depresión muestran un solapamiento considerable, especialmente en la zona central e izquierda del gráfico, reflejando la proximidad conceptual y clínica entre ambos trastornos.

La visualización mediante \ac{t-SNE} refuerza esta misma interpretación de forma más clara. \ac{TCA} vuelve a formar un clúster más definido y separado, mientras que ansiedad y depresión continúan mostrando una fuerte superposición, lo que anticipa una mayor dificultad para diferenciarlas automáticamente. La clase control aparece distribuida entre distintas regiones del espacio, sin formar un clúster propio claramente delimitado, lo cual resulta coherente al tratarse de una categoría heterogénea que no responde a un único patrón lingüístico específico.

Este análisis confirma que la tarea de clasificación presenta una dificultad inherente derivada del propio solapamiento semántico entre clases, especialmente entre ansiedad y depresión. Asimismo, justifica que ciertos errores de clasificación observados posteriormente no dependan únicamente del modelo utilizado, sino también de la ambigüedad natural presente en los datos.


\section{Entrenamiento y evaluación de modelos BERT}

Tras haber analizado la estructura del dataset, identificado sus principales limitaciones y dificultades, y construido los subconjuntos específicos para cada tarea de clasificación, el siguiente paso consiste en aplicar el proceso de entrenamiento y evaluación de los modelos \ac{BERT} seleccionados. En esta fase, los textos preparados se transforman en representaciones adecuadas para el modelo mediante tokenización, se ajustan los modelos preentrenados mediante fine-tuning supervisado y, finalmente, se generan predicciones sobre los conjuntos de prueba para evaluar su rendimiento.

\subsection{Configuración experimental}

El entrenamiento de los modelos \ac{BERT} se realizó utilizando Google Colab como entorno de ejecución, empleando una \ac{GPU} NVIDIA T4. Esta elección permitió disponer de aceleración por hardware suficiente para llevar a cabo el proceso de fine-tuning de modelos Transformer de tamaño medio y grande, aunque con ciertas limitaciones de memoria que condicionaron parte de la configuración experimental.

Tal y como se explicó en la Sección~\ref{subsec:modelos_bert}, en este trabajo se utilizaron cuatro arquitecturas representativas: \ac{DistilBETO}, \ac{RoBERTa}, \ac{mDeBERTa} y Longformer. Aunque todas siguen un enfoque \textit{encoder-only} orientado a tareas de clasificación, presentan diferencias importantes en tamaño, complejidad y capacidad de procesamiento de secuencias largas, lo que obligó a adaptar algunos hiperparámetros de entrenamiento para cada caso.

La configuración específica de cada modelo se definió principalmente en función de tres factores: el número de parámetros de la arquitectura, la longitud máxima de secuencia que podía procesar y las restricciones de memoria de la \ac{GPU} disponible. \ac{DistilBETO} al ser un modelo más ligero, permitió utilizar tamaños de batch mayores tanto en entrenamiento como en evaluación. En cambio, modelos más pesados como \ac{RoBERTa} o \ac{mDeBERTa} requirieron reducir este valor para evitar errores de memoria.

En el caso de Longformer, diseñado específicamente para secuencias largas, se priorizó aumentar la longitud máxima de entrada hasta 2048 tokens en lugar de los 512 habituales en BERT estándar, lo que obligó a reducir considerablemente el tamaño de batch debido al mayor consumo de memoria. Para los cuatro modelos, durante la fase de evaluación se utilizó un batch superior al de entrenamiento, ya que al no calcular gradientes el consumo de memoria resulta menor.

La Tabla~\ref{tab:bert_training_config} resume la configuración utilizada para cada uno de los modelos evaluados.

\begin{table}[H]
	\centering
	\caption{Configuración experimental utilizada para el fine-tuning de los modelos BERT}
	\label{tab:bert_training_config}
	\begin{tabular}{|l|c|c|c|}
		\hline
		\textbf{Modelo} & \textbf{Max length} & \textbf{Train batch size} & \textbf{Eval batch size} \\
		\hline
		DistilBETO & 512 & 16 & 64 \\
		RoBERTa & 512 & 8 & 32 \\
		mDeBERTa & 512 & 8 & 32 \\
		Longformer & 2048 & 1 & 2 \\
		\hline
	\end{tabular}
\end{table}

\subsection{Preparación para el entrenamiento}

Una vez definida la configuración experimental, el siguiente paso consistió en adaptar los textos del dataset al formato requerido por los modelos Transformer. Para ello, se cargó tanto el tokenizador como el modelo preentrenado correspondiente a cada arquitectura utilizando la librería \texttt{Transformers} de Hugging Face \cite{wolf_transformers_2020}.

\subsubsection{Carga del modelo preentrenado}

Cada modelo se inicializó mediante \texttt{AutoModelForSequenceClassification}, ajustando automáticamente la capa de salida al número de clases de la tarea concreta. En los experimentos de clasificación binaria se utilizaron dos etiquetas (\textit{control} y \textit{suffer}), mientras que en la clasificación multiclase se configuró una salida de cuatro clases (control, ansiedad, depresión y \ac{TCA}).

De este modo, el modelo conserva los pesos aprendidos durante la fase de preentrenamiento general, reutilizando el conocimiento lingüístico previamente adquirido. Por su parte, la capa final de clasificación se adapta durante el fine-tuning a la tarea supervisada específica. Durante este proceso, todos los parámetros del modelo pueden ajustarse, permitiendo especializar la representación interna hacia el problema concreto de clasificación de cada trastorno de salud mental.

\subsubsection{Tokenización y adaptación del dataset}

Posteriormente, se aplicó el proceso de tokenización sobre los textos concatenados de cada usuario. Esta fase transforma el texto en una secuencia de identificadores numéricos (\textit{input IDs}) junto con la información auxiliar necesaria para que el modelo pueda procesarlos correctamente, como las máscaras de atención (\textit{attention masks}).

Durante la tokenización se aplicó truncado (\textit{truncation}) y relleno (\textit{padding}) hasta la longitud máxima fija (\textit{max\_length}) definida según la configuración de cada modelo.

Una vez tokenizados los textos, fue necesario adaptar la estructura del dataset al formato esperado por el \texttt{Trainer} de Hugging Face. En clasificación binaria, la columna \texttt{bs} se renombró como \texttt{labels}, mientras que en clasificación multiclase se renombró la columna \texttt{label\_mc} con el mismo objetivo. Además, se eliminaron columnas auxiliares que no eran necesarias para el entrenamiento, como \texttt{texto}, \texttt{user\_id} o \texttt{trastorno}, conservando únicamente la información relevante para la inferencia supervisada.

Finalmente, los datos se convirtieron al formato tensorial de PyTorch \cite{paszke_pytorch_2019} mediante \texttt{set\_format("torch")}, permitiendo su uso directo dentro del proceso de fine-tuning con el \texttt{Trainer} de Hugging Face.

\subsection{Fine-tuning supervisado}

Una vez preparados el modelo y el dataset, se procedió al proceso de fine-tuning supervisado. El objetivo de esta fase es especializar la representación lingüística general aprendida durante el preentrenamiento hacia la identificación de patrones asociados a ansiedad, depresión y \ac{TCA}.

\subsubsection{Ponderación de clases}

Como se observó durante el análisis del dataset, la distribución entre clases \textit{control} y \textit{suffer} no es equilibrada, lo que puede afectar negativamente al aprendizaje supervisado. Si se utiliza una función de pérdida estándar, el modelo tiende a optimizar principalmente el rendimiento sobre la clase más frecuente, reduciendo la sensibilidad sobre las clases minoritarias.

% Formula?
Para mitigar este problema, se calcularon pesos de clase inversamente proporcionales a la frecuencia de cada etiqueta en el conjunto de entrenamiento. De este modo, los errores cometidos sobre clases menos representadas reciben una penalización mayor durante el entrenamiento, obligando al modelo a prestar mayor atención a estos casos.

Estos pesos se incorporaron posteriormente a la función de pérdida mediante una versión ponderada de \texttt{CrossEntropyLoss}, integrada dentro de un \texttt{Trainer} personalizado de Hugging Face denominado \texttt{WeightedTrainer}. Esta modificación permitió mantener el flujo estándar de entrenamiento, adaptando únicamente el cálculo de la pérdida.

\subsubsection{Configuración del entrenamiento}

El entrenamiento se configuró utilizando la clase \texttt{TrainingArguments} de Hugging Face, definiendo una estrategia de evaluación y guardado por épocas (\textit{eval\_strategy="epoch"} y \textit{save\_strategy="epoch"}). De esta forma, al finalizar cada época se evaluaba el modelo sobre el conjunto de validación y se almacenaba un punto de control (\textit{checkpoint}) del entrenamiento.

Los principales hiperparámetros utilizados durante el entrenamiento fueron los siguientes:

\begin{itemize}
	\item \textbf{Número de épocas:} se estableció un máximo de 10 épocas, es decir, 10 recorridos completos sobre el conjunto de entrenamiento. En cada época el modelo ajusta sus parámetros a partir de los ejemplos disponibles, por lo que un número excesivo puede provocar sobreajuste, mientras que un número demasiado bajo puede impedir una correcta convergencia.
	
	\item \textbf{Tasa de aprendizaje (\textit{learning rate}):} se fijó en $2 \times 10^{-5}$. Este parámetro controla cuánto se modifican los pesos del modelo en cada paso de optimización. Valores demasiado altos pueden generar un entrenamiento inestable, mientras que valores demasiado bajos pueden ralentizar excesivamente el aprendizaje o impedir alcanzar un buen mínimo.
	
	\item \textbf{Warmup ratio:} se utilizó un 10\% del entrenamiento como fase de calentamiento. Durante este periodo inicial, la tasa de aprendizaje aumenta progresivamente desde un valor muy bajo hasta alcanzar su valor objetivo. Esto ayuda a estabilizar las primeras actualizaciones del modelo y evita cambios bruscos en los pesos al inicio del fine-tuning.
	
	\item \textbf{Weight decay:} se configuró con un valor de 0.01 como mecanismo de regularización. Su función consiste en penalizar pesos excesivamente grandes durante el entrenamiento, favoreciendo soluciones más estables y reduciendo el riesgo de sobreajuste sobre el conjunto de entrenamiento.
	
	\item \textbf{Precisión mixta (\textit{fp16}):} se activó el entrenamiento en media precisión utilizando representaciones numéricas de 16 bits en lugar de 32 bits para parte de los cálculos. Esto permite reducir el consumo de memoria de la \ac{GPU} y acelerar el entrenamiento, algo especialmente importante al trabajar con modelos grandes en entornos con recursos limitados como Google Colab.
\end{itemize}

Como métrica principal de selección se utilizó el F1-score macro sobre el conjunto de validación, configurando \texttt{load\_best\_model\_at\_end=True} para recuperar automáticamente el mejor checkpoint al finalizar el entrenamiento. Esta elección resulta especialmente adecuada en escenarios con desbalance de clases, ya que evita que el modelo optimice únicamente la precisión (\textit{accuracy}) global.

\subsubsection{Early stopping y control del sobreajuste}

Para evitar entrenamientos innecesariamente largos y reducir el riesgo de sobreajuste, se incorporó un mecanismo de \textit{early stopping} con una paciencia de dos épocas. Esto significa que, si el valor de F1-score sobre validación no mejoraba durante dos evaluaciones consecutivas, el entrenamiento se detenía automáticamente.

Esta estrategia permitió estabilizar el proceso de fine-tuning y asegurar que el modelo final seleccionado correspondiera al mejor rendimiento observado durante el entrenamiento, en lugar de utilizar simplemente la última época ejecutada.

\subsection{Evaluación e inferencia sobre el conjunto de prueba}

Una vez finalizado el proceso de fine-tuning, el modelo seleccionado se evaluó sobre el conjunto de prueba (\textit{test set}), que no había sido utilizado ni durante el entrenamiento ni durante la validación. Esta fase permite estimar el rendimiento real del modelo sobre datos no vistos previamente y obtener las predicciones finales que posteriormente serán analizadas en el capítulo de resultados.

La inferencia se realizó utilizando el método \texttt{predict()} del \texttt{Trainer} de Hugging Face, que genera automáticamente los logits asociados a cada ejemplo del conjunto de prueba. Estos logits representan las puntuaciones sin normalizar asignadas a cada clase posible por el modelo.

Posteriormente, las predicciones finales se obtuvieron aplicando la operación \texttt{argmax} sobre dichos logits, seleccionando para cada ejemplo la clase con mayor puntuación. En clasificación binaria, esto permitió asignar cada usuario a la clase \textit{control} o \textit{suffer}, mientras que en clasificación multiclase se seleccionó entre control, ansiedad, depresión y \ac{TCA}.

Las predicciones generadas se incorporaron al dataset original de prueba mediante una nueva columna denominada \texttt{pred}, conservando así tanto las etiquetas reales como las predicciones del modelo en un mismo fichero. Posteriormente, estos resultados se almacenaron en formato CSV para su análisis posterior y comparación con los resultados obtenidos mediante \acp{LLM}.

\subsection{Tiempos de entrenamiento e inferencia}

Además del rendimiento predictivo, durante el proceso de fine-tuning se registraron los tiempos de entrenamiento e inferencia de cada modelo con el objetivo de analizar el coste práctico asociado a cada arquitectura. Esta información resulta especialmente relevante al comparar modelos con diferentes tamaños y capacidades de procesamiento, ya que una mejora en las métricas de clasificación no siempre compensa un incremento significativo en tiempo de ejecución.

El tiempo de entrenamiento corresponde al proceso completo de fine-tuning sobre el conjunto de entrenamiento y validación, incluyendo la selección automática del mejor modelo mediante \textit{early stopping}. Por su parte, el tiempo de inferencia se midió durante la generación de predicciones sobre el conjunto de prueba utilizando el método \texttt{predict()} del \texttt{Trainer}. La suma de ambos valores permite obtener una estimación del coste temporal total de cada experimento.

La Tabla~\ref{tab:bert_times} recoge los tiempos registrados para cada combinación de modelo y tarea de clasificación.

\begin{table}[H]
	\centering
	\caption{Tiempos de entrenamiento e inferencia de los modelos BERT}
	\label{tab:bert_times}
	\begin{tabular}{|l|c|c|c|}
		\hline
		\textbf{Modelo} & \textbf{Tiempo Entr. (min)} & \textbf{Tiempo Eval. (min)} & \textbf{Tiempo Total (min)} \\
		\hline
		
		DistilBETO - Anxiety &  &  &  \\
		DistilBETO - Depression &  &  &  \\
		DistilBETO - ED &  &  &  \\
		DistilBETO - Multiclass &  &  &  \\
		\hline
		
		RoBERTa - Anxiety &  &  &  \\
		RoBERTa - Depression &  &  &  \\
		RoBERTa - ED &  &  &  \\
		RoBERTa - Multiclass &  &  &  \\
		\hline
		
		mDeBERTa - Anxiety &  &  &  \\
		mDeBERTa - Depression &  &  &  \\
		mDeBERTa - ED &  &  &  \\
		mDeBERTa - Multiclass &  &  &  \\
		\hline
		
		Longformer - Anxiety &  &  &  \\
		Longformer - Depression &  &  &  \\
		Longformer - ED &  &  &  \\
		Longformer - Multiclass &  &  &  \\
		\hline
		
	\end{tabular}
\end{table}


\section{Inferencia con LLMs mediante prompting}


%%% Local Variables:
%%% TeX-master: "../book"
%%% End:
